\section{用 Fourier 变换构造}\srcpage{11}

对每个 $i=1,\dots,n$，设 $\Coh_i(X)$ 为一般秩为 $i$ 的凝聚 $\cO_X$-模 $\cM_i$ 构成的代数叠。
设 ${\Coh'}_i(X)$ 为一般秩为 $i$ 的凝聚 $\cO_X$-模 $\cM_i$ 再附上一个模数乘的
$\cO_X$-线性单射 $s_i\colon(\OmX)^{\otimes i-1}\hookrightarrow\cM_i$ 所得的代数叠。
这两个叠以 $\Fib_i(X)$ 与 ${\Fib'}_i(X)$ 为稠密开集。
与后两个叠一样，它们的连通分支 $\Coh^d_i(X)$ 与 ${\Coh'}^d_i(X)$ 按 $\cM_i$ 的次数 $d$ 划分。

\srcpage{12}
\begin{lemma}
存在常数 $c(g,n)$，使得以下结论成立。
对每个整数 $d\ge c(g,n)$ 与每个 $L\in\Fib^d_n(X)$，只要
$\Hom_{\cO_X}\big(L,(\OmX)^{\otimes n+1}\big)\ne(0)$，$L$ 就是非常不稳定的。
也就是说，存在非平凡的直和分解 $L=L_1\oplus L_2$，使 $\Ext^1_{\cO_X}(L_1,L_2)=(0)$。%
\hfill$\square$
\end{lemma}

对每个 $i=0,\dots,n$，设 $\cC_i\subset\Coh_i(X)$ 为由满足以下两条的 $\cM_i$ 构成的开集：
其次数 $\ge c(n,g)+i(i-1)(g-1)$，且 $\Hom_{\cO_X}\big(\cM_i,(\OmX)^{\otimes n+1}\big)=(0)$。
于是特别地 $\Hom_{\cO_X}\big(\cM_i,(\OmX)^{\otimes i}\big)=(0)$；
并且由 Serre 对偶，$\Ext^1_{\cO_X}\big((\OmX)^{\otimes i-1},\cM_i\big)=(0)$。
设 $\cV_i$ 为二元组 $(\cM_i,s_i)$ 的代数叠，其中 $\cM_i\in\cC_i$，而
$s_i\in\Hom_{\cO_X}\big((\OmX)^{\otimes i-1},\cM_i\big)$。
设 $\cV^\vee_i$ 为满足 $\cM_i\in\cC_i$ 的凝聚 $\cO_X$-模扩张
\[
  0\to(\OmX)^{\otimes i}\to\cM_{i+1}\to\cM_i\to 0
\]
的代数叠。有自然的投影 $\pi_i\colon\cV_i\to\cC_i$ 与 $\pi^\vee_i\colon\cV^\vee_i\to\cC_i$，
它们是对偶的向量丛。

引入开集
\[
  \cV{}^\circ_i=\big\{(\cM_i,s_i)\ \big|\ s_i\ \text{是单射}\big\}\subset\cV_i
\]
与
\[
  \cV{}^{\vee\circ}_i=\big\{(0\to(\OmX)^{\otimes i}\to\cM_{i+1}\to\cM_i\to 0)\ \big|\
  \cM_{i+1}\in\cC_{i+1}\big\}\subset\cV^\vee_i .
\]
当然，有典范同构
\[
  \cV{}^{\vee\circ}_i\xrightarrow{\ \sim\ }\cV{}^\circ_{i+1},
\]
它把 $0\to(\OmX)^{\otimes i}\to\cM_{i+1}\to\cM_i\to 0$ 送到单射
$(\OmX)^{\otimes i}\to\cM_{i+1}$。

乘法群按数乘作用在向量丛 $\cV_i$ 与 $\cV^\vee_i$ 上；取商即得基本图
\begin{center}
\footnotesize
\begin{tikzcd}[column sep=1.55em, row sep=1.7em]
  & \bar\cV{}^\circ_n
      \arrow[dl]
  & \bar\cV{}^{\vee\circ}_{n-1}
      \arrow[l, "\sim"']
      \arrow[r, hook]
  & \bar\cV{}^\vee_{n-1}
      \arrow[dr]
  & & \bar\cV_{n-1}
      \arrow[dl]
  & \bar\cV{}^\circ_{n-1}
      \arrow[l, hook']
  & \bar\cV{}^{\vee\circ}_{n-2}
      \arrow[l, "\sim"']
      \arrow[r, hook]
  & \cdots
  & & & \\
  \cC_n & & & & \cC_{n-1} & & & & & & & \\
  \cdots
  & \bar\cV{}^\vee_2
      \arrow[dr]
  & & \bar\cV_2
      \arrow[dl]
  & \bar\cV{}^\circ_2
      \arrow[l, hook']
  & \bar\cV{}^{\vee\circ}_1
      \arrow[l, "\sim"']
      \arrow[r, hook]
  & \bar\cV{}^\vee_1
      \arrow[dr]
  & & \bar\cV_1
      \arrow[dl]
  & \bar\cV{}^\circ_1
      \arrow[l, hook']
  & \bar\cV{}^{\vee\circ}_0
      \arrow[l, "\sim"']
      \arrow[dr]
  & \\
  & & \cC_2 & & & & & \cC_1 & & & & \cC_0
\end{tikzcd}
\end{center}
其中斜箭头是投影
$\bar\pi{}^\circ_n\colon\bar\cV{}^\circ_n=[\cV^\circ_n/\Gm]\to\cC_n$，
$\bar\pi_i\colon\bar\cV_i=[\cV_i/\Gm]\to\cC_i$，
$\bar\pi{}^\vee_i\colon\bar\cV{}^\vee_i=[\cV^\vee_i/\Gm]\to\cC_i$，
$\bar\pi{}^{\vee\circ}_0\colon\bar\cV{}^{\vee\circ}_0=[\cV^{\vee\circ}_0/\Gm]\to\cC_0$。
水平箭头一方面是开浸入
$j_i\colon\bar\cV{}^\circ_i\hookrightarrow\bar\cV_i$ 与
$j^\vee_i\colon\bar\cV{}^{\vee\circ}_i\hookrightarrow\bar\cV{}^\vee_i$，
它们由包含 $\cV^\circ_i\subset\cV_i$ 与 $\cV^{\vee\circ}_i\subset\cV^\vee_i$ 诱导；
另一方面是同构
$\iota_i\colon\bar\cV{}^{\vee\circ}_i\xrightarrow{\ \sim\ }\bar\cV{}^\circ_{i+1}$，
它们由上面的同构 $\cV{}^{\vee\circ}_i\xrightarrow{\ \sim\ }\cV{}^\circ_{i+1}$ 诱导。

\srcpage{13}
于是对 $X$ 上每个秩为 $n$ 的局部系统 $E$，定义复形
$K_{E,i}\in\Dbc(\bar\cV{}^\circ_i)$，$i=1,\dots,n$，如下：
\[
  K_{E,1}=(\iota_0)_*\,(\bar\pi{}^{\vee\circ}_0)^*\,\cL_E,
\]
其中 $\cL_E$ 是 $\cC_0$ 上的复形：它在 $\cC_0$ 的每个连通分支 $\cC^d_0$ 上的限制，
是 3.2 中定义的反常层 $\cL^d_E$ 平移 $d-1$；并由递推关系
\[
  K_{E,i+1}=(\iota_i)_*\,(j^\vee_i)^*\,\Four_{\bar\cV_i/\cC_i}\big((j_i)_!K_{E,i}\big),
  \qquad i=1,\dots,n-1
\]
定义，其中 $\Four_{\bar\cV_i/\cC_i}\colon\Dbc(\bar\cV_i)\to\Dbc(\bar\cV{}^\vee_i)$ 是向量丛
$\cV_i\to\cC_i$ 的齐次 Fourier 变换，其定义见第~0~节。

叠 $\bar\cV{}^\circ_i$ 与 ${\Fib'}_i(X)$ 都是 ${\Coh'}_i(X)$ 的开集，因此可以考虑它们的交。

\begin{lemma}
设 $E$ 是 $X$ 上秩为 $n$ 的局部系统。
对每个 $i=1,\dots,n$ 与每个整数 $d\ge 0$，$K_{E,i}$ 在
${\Coh'}^{i(i-1)(g-1)+d}_i(X)$ 的开集
$\bar\cV{}^\circ_i\cap{\Fib'}^{i(i-1)(g-1)+d}_i(X)$ 上的限制等于
\[
  M'^{d}_{i,E}(L_i),
\]
这里采用第 4 节的记号。

特别地，对每个整数 $d\ge 0$，在 ${\Coh'}^{n(n-1)(g-1)+d}_i(X)$ 的开集
$\bar\cV{}^\circ_n\cap{\Fib'}^{n(n-1)(g-1)+d}_i(X)$ 上，$K_{E,n}$ 与 $\Aut'{}_E$ 的限制相合。%
\yizhu{末句两处下标写的都是 $i$，与句首的 $K_{E,n}$、$\Aut'{}_E$ 以及本引理首段不合，
似应为 $n$；此处照原文翻译。}{原文笔误，I（据 p13 上下文，无外部依据）}\hfill$\square$
\end{lemma}
